\cotefile{62}
\scanpage{51}
\mspage{- 27 -}

n\textsuperscript{o} 276

théorème sur les anneaux de Zariski, plutôt que des disperser les diverses étapes de la démonstration en Propositions auxilliaires ainsi qu'il est fait dans l'Etat Serre ; on peut espérer que le lecteur moyen sera bien capable de voir tout seul, par exemple, que la relation $\hat{\mathrm E}=\overset{\handwritten{\scriptscriptstyle\wedge}}{\mathrm A}\handwritten{\otimes}\mathrm E$ est vraie pour tout anneau $\mathrm m$-adique noethérien, vu que la démonstration ne fait pas usage de l'hypothèse $\mathrm m\handwritten{\subset}\mathrm r(\mathrm A)$ \ldots)

% archive sheet number: 50
